\documentclass[11pt]{article}
% ===========================================================================
%  Shared preamble for the AP Calculus AB exam-preparation worksheets.
%  Input right after \documentclass. Files beginning with "_" are not compiled
%  on their own. Navy and gold, to match the course's lesson pages.
% ===========================================================================
\usepackage[letterpaper,top=2.3cm,bottom=2.2cm,left=1.7cm,right=1.7cm,headheight=28pt,footskip=30pt]{geometry}
\usepackage{amsmath,amssymb}
\usepackage[T1]{fontenc}
\usepackage[utf8]{inputenc}
\usepackage{lmodern}
\usepackage{xcolor}
\usepackage[most]{tcolorbox}
\usepackage{enumitem}
\usepackage{fancyhdr}
\usepackage{array}
\usepackage{tikz}
\usepackage{pgfplots}
\pgfplotsset{compat=1.18}

\definecolor{navy}{HTML}{161E54}
\definecolor{ink}{HTML}{27357E}
\definecolor{gold}{HTML}{B98A00}
\definecolor{goldbg}{HTML}{FFFDF3}
\definecolor{inkbg}{HTML}{F6F7FD}
\definecolor{rule}{HTML}{C9CFF0}
\definecolor{okgreen}{HTML}{1B7A44}
\definecolor{okbg}{HTML}{F1F8F3}

\pagestyle{fancy}
\fancyhf{}
\fancyhead[L]{\footnotesize Name: \underline{\hspace{4cm}}}
\fancyhead[C]{\textbf{\color{navy}Integration Academy}\\[-2pt]{\scriptsize Calculus AB \textperiodcentered{} AP\textsuperscript{\textregistered} Exam Preparation}}
\fancyfoot[L]{\tiny\color{black!55}\parbox[t]{0.78\textwidth}{AP\textsuperscript{\textregistered} is a registered trademark of the College Board, which was not involved in the production of, and does not endorse, this material. \textcopyright{} Integration Academy.}}
\fancyfoot[R]{\scriptsize\thepage}
\renewcommand{\headrulewidth}{1.2pt}
\renewcommand{\headrule}{\hbox to\headwidth{\color{gold}\leaders\hrule height \headrulewidth\hfill}}
\renewcommand{\footrulewidth}{0pt}
\setlength{\parindent}{0pt}
\setlength{\parskip}{4pt}

% Title block: \wstitle{code}{title}{unit line}
\newcommand{\wstitle}[3]{\begin{center}{\small\bfseries\color{gold}WORKSHEET #1}\\[3pt]{\LARGE\bfseries\color{navy}#2}\\[3pt]{\small\color{black!60}#3}\end{center}\vspace{2pt}}
% Section heading
\newcommand{\sech}[1]{\par\vspace{10pt}{\large\bfseries\color{navy}#1}\par\nobreak\vspace{2pt}{\color{rule}\hrule height0.8pt}\vspace{7pt}}
% A level tag printed beside a group of questions
\newcommand{\level}[1]{\par\vspace{6pt}{\footnotesize\bfseries\color{gold}\MakeUppercase{#1}}\par\nobreak\vspace{2pt}}

\newtcolorbox{ideabox}{colframe=ink,colback=inkbg,coltitle=white,colbacktitle=ink,fonttitle=\bfseries\small,title={Key ideas},arc=2pt,boxrule=0.8pt,left=8pt,right=8pt}
\newtcolorbox{example}[1]{colframe=rule,colback=white,coltitle=navy,colbacktitle=inkbg,fonttitle=\bfseries\small,title={#1},arc=2pt,boxrule=0.8pt,borderline west={3pt}{0pt}{ink},left=10pt,right=8pt}
\newtcolorbox{frbox}[1]{colframe=gold,colback=goldbg,coltitle=white,colbacktitle=gold,fonttitle=\bfseries\small,title={#1},arc=2pt,boxrule=0.8pt,breakable,left=8pt,right=8pt}
\newtcolorbox{solbox}{colframe=okgreen,colback=okbg,coltitle=white,colbacktitle=okgreen,fonttitle=\bfseries\small,title={Worked solutions},arc=2pt,boxrule=0.8pt,breakable,left=8pt,right=8pt}
\newtcolorbox{scorebox}[1]{colframe=navy,colback=white,coltitle=white,colbacktitle=navy,fonttitle=\bfseries\small,title={#1},arc=2pt,boxrule=0.8pt,breakable,left=8pt,right=8pt}

\newcommand{\soln}{\par\smallskip\textbf{\color{ink}Solution.}\ }
\newcommand{\result}[1]{\par\smallskip\textbf{#1}}
% blank working space of a given height
\newcommand{\work}[1]{\par\vspace{#1}}
% a point value in the margin of a scoring guide
\newcommand{\pts}[1]{\hfill{\footnotesize\bfseries\color{gold}[#1]}}

% Coordinate grid for graph questions: \begin{apgraph}{xmin}{xmax}{ymin}{ymax} ... \end{apgraph}
\newenvironment{apgraph}[4]{\begin{center}\begin{tikzpicture}\begin{axis}[
  axis lines=middle, xmin=#1, xmax=#2, ymin=#3, ymax=#4,
  grid=both, grid style={black!12}, xtick distance=1, ytick distance=1,
  ticklabel style={font=\scriptsize}, xlabel={$x$}, ylabel={$y$},
  width=10cm, height=7.2cm, samples=80, clip=false,
  every axis plot/.append style={line width=1.1pt, color=ink}]}{\end{axis}\end{tikzpicture}\end{center}}
\newcommand{\closeddot}[2]{\addplot[only marks, mark=*, mark size=2.2pt, color=ink] coordinates {(#1,#2)};}
\newcommand{\opendot}[2]{\addplot[only marks, mark=*, mark size=2.2pt, color=ink, mark options={fill=white}] coordinates {(#1,#2)};}

% Multiple-choice options: four across, or stacked when the options are long.
\newcommand{\opts}[4]{\par\vspace{3pt}\begin{tabular}{@{}p{0.235\linewidth}p{0.235\linewidth}p{0.235\linewidth}p{0.235\linewidth}@{}}(A)\ #1 & (B)\ #2 & (C)\ #3 & (D)\ #4\end{tabular}\par\vspace{9pt}}
\newcommand{\optsv}[4]{\par\vspace{3pt}\begin{tabular}{@{}l@{\ \ }p{0.9\linewidth}@{}}(A) & #1\\(B) & #2\\(C) & #3\\(D) & #4\end{tabular}\par\vspace{9pt}}

\fancyhead[R]{\footnotesize\color{black!60}Worksheet 8.7}
\begin{document}

\wstitle{8.7}{Volumes from Known Cross-Sections}{Unit 8: Applications of Integration}

\begin{ideabox}
\begin{itemize}[leftmargin=*,itemsep=2pt,topsep=2pt]
  \item If the cross-section at $x$ has area $A(x)$, the volume is $\displaystyle\int_{a}^{b}A(x)\,dx$.
  \item The side $s$ of each cross-section lies in the base: $s=\text{top}-\text{bottom}$.
  \item Square: $s^2$. \quad Semicircle of diameter $s$: $\dfrac{\pi}{8}s^2$. \quad Equilateral triangle: $\dfrac{\sqrt3}{4}s^2$.
  \item Isosceles right triangle with a leg in the base: $\dfrac12s^2$. \quad Rectangle of height $h$: $sh$.
\end{itemize}
\end{ideabox}

\sech{Worked examples}
\begin{example}{Example 1 \textperiodcentered{} Squares}
The base of a solid is bounded by $y=x+1$, the $x$-axis, $x=0$ and $x=3$. Cross-sections perpendicular to the $x$-axis are squares. Find the volume.\soln $\displaystyle\int_{0}^{3}(x+1)^2\,dx=\dfrac{64-1}{3}$.\result{$21$.}
\end{example}

\begin{example}{Example 2 \textperiodcentered{} Semicircles}
The base is bounded by $y=\sqrt{x}$, the $x$-axis and $x=9$. Cross-sections perpendicular to the $x$-axis are semicircles with diameter in the base. Find the volume.\soln $\dfrac\pi8\displaystyle\int_{0}^{9}x\,dx=\dfrac\pi8\cdot\dfrac{81}{2}$.\result{$\dfrac{81\pi}{16}$.}
\end{example}

\begin{example}{Example 3 \textperiodcentered{} Perpendicular to the $y$-axis}
The base lies in the first quadrant, bounded by $y=x^2$, the $y$-axis and $y=9$. Cross-sections perpendicular to the $y$-axis are squares. Find the volume.\soln The side is $x=\sqrt y$: $\displaystyle\int_{0}^{9}y\,dy$.\result{$\dfrac{81}{2}$.}
\end{example}

\newpage
\sech{Practice}

\level{Core}
\begin{enumerate}[leftmargin=*,itemsep=2pt]
  \item Find the area of a semicircle with diameter 6.\work{2.3cm}
  \item Find the area of an equilateral triangle with side 4.\work{2.3cm}
  \item A solid has cross-sectional area $A(x)=3x^2$ for $0\le x\le2$. Find its volume.\work{2.3cm}
  \item A solid has square cross-sections of side $2x$ for $0\le x\le3$. Find its volume.\work{2.3cm}
\end{enumerate}

\level{Exam level}
\begin{enumerate}[leftmargin=*,itemsep=2pt,start=5]
  \item The base of a solid is the disc $x^2+y^2\le25$. Cross-sections perpendicular to the $x$-axis are semicircles with diameter in the base. Find the volume.\work{2.8cm}
  \item The base is the region enclosed by $y=x$ and $y=x^2$. Cross-sections perpendicular to the $x$-axis are isosceles right triangles with a leg in the base. Find the volume.\work{2.8cm}
  \item The base is bounded by $y=6-2x$ and the coordinate axes. Cross-sections perpendicular to the $x$-axis are equilateral triangles. Find the volume.\work{2.8cm}
  \item \emph{Calculator.} The base is bounded by $y=\ln x$, the $x$-axis and $x=e$. Cross-sections perpendicular to the $x$-axis are squares. Find the volume to three decimal places.\work{2.8cm}
\end{enumerate}

\begin{frbox}{Question 9 \textperiodcentered{} Free-response practice}
Let $R$ be the region bounded by $y=6x-x^2$ and the $x$-axis. $R$ is the base of a solid.
\begin{enumerate}[label=(\alph*),leftmargin=*,itemsep=2pt]
  \item Find the area of $R$.\work{2.6cm}
  \item Cross-sections of the solid perpendicular to the $x$-axis are squares. Find the volume.\work{2.6cm}
  \item For a different solid with the same base, the cross-section at $x$ is a rectangle whose height is $x$. Find the volume.\work{2.6cm}
\end{enumerate}
\end{frbox}

\newpage
\begin{solbox}
\begin{enumerate}[leftmargin=*,itemsep=4pt]
  \item $\dfrac12\pi(3)^2=\dfrac{9\pi}{2}$.
  \item $\dfrac{\sqrt3}{4}(16)=4\sqrt3$.
  \item $\Big[x^3\Big]_0^2=8$.
  \item $\displaystyle\int_{0}^{3}4x^2\,dx=36$.
  \item Diameter $2\sqrt{25-x^2}$: $\dfrac\pi8\displaystyle\int_{-5}^{5}4(25-x^2)\,dx=\dfrac\pi2\cdot\dfrac{500}{3}=\dfrac{250\pi}{3}$ (half of a sphere of radius 5).
  \item $\dfrac12\displaystyle\int_{0}^{1}(x-x^2)^2\,dx=\dfrac12\left(\dfrac13-\dfrac12+\dfrac15\right)=\dfrac{1}{60}$.
  \item $\dfrac{\sqrt3}{4}\displaystyle\int_{0}^{3}(6-2x)^2\,dx=\dfrac{\sqrt3}{4}(36)=9\sqrt3$.
  \item $\displaystyle\int_{1}^{e}(\ln x)^2\,dx\approx0.718$.
  \item \begin{enumerate}[label=(\alph*),leftmargin=*,itemsep=3pt]
    \item $\displaystyle\int_{0}^{6}(6x-x^2)\,dx=108-72=36$.
    \item $\displaystyle\int_{0}^{6}(36x^2-12x^3+x^4)\,dx=2592-3888+\dfrac{7776}{5}=\dfrac{1296}{5}$.
    \item $\displaystyle\int_{0}^{6}x(6x-x^2)\,dx=432-324=108$.
  \end{enumerate}
\end{enumerate}
\end{solbox}

\end{document}
